% !TeX root = ../navier-stokes-manuscript.tex
\subsection{The Heat Clock and the Zeno Cascade}\label{sub:BodyHeatClock}

The smaller width comes with a shorter viscous clock.
This is already visible in the linear heat equation.
If \(v_t=\nu\Delta v\) and the Fourier mass of \(v\) lies near \(|\xi|\simeq r^{-1}\), then
\[
  \widehat v(\xi,t+\delta t)
  =e^{-\nu|\xi|^2\delta t}\widehat v(\xi,t).
\]
An order-one heat damping occurs when \(\nu|\xi|^2\delta t\) is of order one.
A spatial width \(r\) is accordingly assigned the comparison time
\begin{equation}\label{eq:BodyHeatComparisonTime}
  \tau_\nu(r)\asymp\frac{r^2}{\nu}.
\end{equation}
Halving the width quarters this clock.

That is a comparison clock.
The heat calculation does not create a smaller width in the nonlinear solution, and it does not say how long the solution remains at a width it has formed.
It says that once structure at width \(r\) is present, viscosity responds to that structure on the scale \(r^2/\nu\).
The same narrowing that can make a spatial gradient larger also makes viscous action faster.

Now place geometric widths beside those clocks:
\[
  r_n\asymp2^{-n}r_0,
  \qquad
  \tau_{\nu,n}:=\tau_\nu(r_n).
\]
Then
\[
  \tau_{\nu,n}
  \asymp\frac{r_0^2}{\nu}4^{-n},
  \qquad
  \sum_{n=0}^{\infty}\tau_{\nu,n}<\infty.
\]
The clocks attached to successively smaller widths fit inside a finite total time.
That is the arithmetic behind the Zeno cascade.

\Needspace{15\baselineskip}
The arithmetic has not yet produced a cascade.
An actual cascade would require one velocity history to form widths
\[
  r_0>r_1>r_2>\cdots\longrightarrow0
\]
at times
\[
  t_0<t_1<t_2<\cdots\uparrow\tau<\infty,
\]
with the transition intervals themselves following a summable schedule.
Neither the widths nor the relation between \(t_{n+1}-t_n\) and \(\tau_{\nu,n}\) comes from the heat multiplier alone.
Those are claims about the coupled Navier--Stokes evolution, and every transition would have to occur in the same fluid while viscosity is already acting on the newly formed scale.

The full scaling still shows why this finite-time geometry is dangerous.
Fix a nonzero divergence-free profile \(\phi\in\Ssigma\) whose Fourier support lies in a fixed annulus, and for \(0<\ell<1\) set
\[
  u_\ell(x)
  :=\ell^{-1}\phi\!\left(\frac{x-x_\ell}{\ell}\right).
\]
The annular support keeps each copy on one frequency scale, so \(\ell\) is the width being compared.
This is the spatial face of one normalized profile under the Navier--Stokes rescaling.
It is a family of comparison fields, not a sequence already realized by the fixed solution.
Its scales are exact:
\[
\begin{aligned}
  \norm{u_\ell}_2^2
  &=\ell\norm{\phi}_2^2,
  &
  \norm{u_\ell}_{\dot H^{1/2}}^2
  &=\norm{\phi}_{\dot H^{1/2}}^2,
  \\
  \norm{u_\ell}_{L^\infty}
  &=\ell^{-1}\norm{\phi}_{L^\infty},
  &
  \norm{\nabla u_\ell}_{L^\infty}
  &=\ell^{-2}\norm{\nabla\phi}_{L^\infty}.
\end{aligned}
\]
The energy of the profile decreases with \(\ell\), its critical height stays fixed, its velocity amplitude grows like \(\ell^{-1}\), and its gradient grows like \(\ell^{-2}\).
Its viscous comparison time is \(\ell^2/\nu\).
The same scale change has made the local field larger, its spatial variation sharper, and the available clock shorter, even though the kinetic energy carried by the profile is disappearing.

This normalized family displays the direction of concentration.
Its kinetic energy tends to zero while its amplitude and gradient diverge.
It does not say that the fixed Navier--Stokes history realizes this sequence, and its fixed critical height is not the growth of \(H(t)\) along an actual history.
It says why the energy law cannot rule such concentration out.
